% =====================================================================
%  DeepKnight technical report, adapted from the final group-project report.
%  Historical measurements refer to the July 2026 implementation.
%  Compile: tectonic technical-report.tex (or pdflatex twice).
% =====================================================================
\documentclass[11pt,a4paper]{article}

\usepackage[utf8]{inputenc}
\usepackage[T1]{fontenc}
\usepackage[ngerman]{babel}
\usepackage{lmodern}
\usepackage[a4paper,margin=2.5cm]{geometry}
\usepackage{booktabs}
\usepackage{array}
\usepackage{amsmath}
\usepackage{amssymb}
\usepackage{listings}
\usepackage{xcolor}
\usepackage{enumitem}
\usepackage{fancyhdr}
\usepackage{caption}
\usepackage{graphicx}
\usepackage{tikz}
\usetikzlibrary{arrows.meta,positioning,fit,backgrounds,calc}
\usepackage[hidelinks]{hyperref}

% ---- Kopf-/Fusszeile -------------------------------------------------
\pagestyle{fancy}
\fancyhf{}
\lhead{\small DeepKnight}
\rhead{\small Technischer Bericht}
\cfoot{\thepage}
\renewcommand{\headrulewidth}{0.4pt}

% ---- Code-Listing-Stil ----------------------------------------------
\definecolor{codegray}{rgb}{0.95,0.95,0.95}
\definecolor{codekw}{rgb}{0.0,0.2,0.6}
\definecolor{codecmt}{rgb}{0.30,0.55,0.30}
\lstset{
  basicstyle=\ttfamily\footnotesize,
  backgroundcolor=\color{codegray},
  keywordstyle=\color{codekw}\bfseries,
  commentstyle=\color{codecmt}\itshape,
  breaklines=true,
  frame=single,
  framesep=4pt,
  rulecolor=\color{gray!40},
  showstringspaces=false,
  language=Python,
  columns=fullflexible,
  keepspaces=true,
}

\newcommand{\code}[1]{\texttt{\small #1}}

% Lange \texttt{}-Bezeichner (z.B. Funktionsnamen) brechen nicht um; diese
% Feineinstellungen lassen TeX enge Zeilen leicht dehnen statt \"uber den Rand
% laufen zu lassen -- vermeidet Overfull-\hbox ohne sichtbar lockere Absaetze.
\setlength{\emergencystretch}{3em}
\tolerance=1200

% =====================================================================
\begin{document}

% ---- Titelseite ------------------------------------------------------
\begin{titlepage}
\thispagestyle{empty}
\centering
\vspace*{1.5cm}

{\large\textsc{DeepKnight}}\\[1.5cm]

\rule{\linewidth}{0.4pt}\\[0.4cm]
{\normalsize Tablut-KI mit Bitboards, PVS und MCTS/RAVE}\\[0.4cm]
\rule{\linewidth}{0.4pt}\\[2.5cm]

{\Huge\bfseries Technischer Bericht}\\[0.8cm]
{\Large Entwicklung einer spielstarken KI}\\[0.2cm]
{\Large f\"ur das Spiel Tablut}\\[0.4cm]
{\large\bfseries Architektur, Suchverfahren und Evaluation}\\[2.0cm]

{\normalsize Portfolio-Ausgabe von Yousef Mousa}\\[0.5cm]
{\small Basierend auf dem gemeinsamen TU-Berlin-Projekt mit}\\[0.2cm]
{\small Batuhan Caniberk und Adham Ahmed Mostafa Mohamed Ibrahim}

\vfill
{\normalsize Repository: \href{https://github.com/Robertsvg/deepknight}{github.com/Robertsvg/deepknight}}\\[0.3cm]
{\small Technische Ergebnisse der Projektfassung vom Juli 2026.}\\[0.2cm]
{\small Aktuelle Regressionstests und Wartungs\"anderungen: siehe README.}
\end{titlepage}

% ---- Zusammenfassung / Abstract -------------------------------------
\begin{abstract}
\noindent
Dieser Bericht dokumentiert die Entwicklung von \textbf{DeepKnight}, einer
spielstarken KI f\"ur das asymmetrische Wikinger-Brettspiel \emph{Tablut}, \"uber
vier Meilensteine hinweg.
Der Ansatz ist bewusst \emph{symbolisch} -- Spielbaumsuche mit handgeschriebener
Bewertungsfunktion, ohne maschinelles Lernen.
Kern ist ein Alpha-Beta-Stack mit iterativer Tiefensuche, Transposition Table,
Principal Variation Search sowie erweiterter Zugsortierung (Killer Moves und
History-Heuristik), erg\"anzt um ein spielphasenabh\"angiges Zeitmanagement.
Die Zustandsrepr\"asentation nutzt \textbf{Bitboards} (drei 81-Bit-Integer), was
die Bewertung gegen\"uber der objektbasierten Fassung um rund das
\textbf{72-Fache} beschleunigt.
Im vierten Meilenstein kamen eine eigenst\"andige Anbindung an den offiziellen
TU-Gameserver mit \emph{server-treuer} Schlaglogik (\textbf{0 Desyncs} \"uber
$152\,000$ verifizierte Z\"uge), eine im Contest geh\"artete Wei\ss-Verteidigung
(die \glqq K\"onig-Falle\grqq{}) sowie eine komplement\"are Monte-Carlo-Baumsuche
mit RAVE-Erweiterung hinzu.
Die finale Projektfassung ist durch \textbf{138 automatisierte Tests} abgesichert.
In direkten Contests gegen eine andere Gruppe gewann DeepKnight nach der
Eval-H\"artung in beiden Farbrollen.
\end{abstract}

% ---- Inhaltsverzeichnis ---------------------------------------------
\tableofcontents
\newpage

% =====================================================================
\section{Einf\"uhrung}
% =====================================================================
Im Rahmen des Projekts \emph{Symbolische K\"unstliche Intelligenz} hat Gruppe~A
\"uber vier Meilensteine hinweg eine spielstarke KI f\"ur das Wikinger-Brettspiel
\textbf{Tablut} entwickelt.
Es handelt sich bewusst um \emph{klassische, symbolische} KI ohne maschinelles
Lernen: Spielbaumsuche mit einer handgeschriebenen Bewertungsfunktion.
Die KI tr\"agt den Namen \textbf{DeepKnight} und ist vollst\"andig in Python~3
implementiert.

Der Bericht fasst den gesamten Projektverlauf zusammen und ist wie folgt
gegliedert.
Abschnitt~\ref{sec:tablut} beschreibt die Spielregeln.
Abschnitt~\ref{sec:arch} erl\"autert die Zustandsrepr\"asentation und die
Systemarchitektur.
Abschnitt~\ref{sec:search} stellt die Suchverfahren dar, die \"uber die
Meilensteine hinweg entstanden sind (Minimax, Alpha-Beta, Transposition Table,
Principal Variation Search, erweiterte Zugsortierung, Spielphasen-Zeitmanagement).
Abschnitt~\ref{sec:eval} dokumentiert die Entwicklung der Bewertungsfunktion.
Abschnitt~\ref{sec:contest} beschreibt die im vierten Meilenstein neu
entwickelte Contest-Anbindung an den offiziellen TU-Gameserver.
Abschnitt~\ref{sec:mcts} behandelt die komplement\"are Technik
Monte-Carlo-Baumsuche.
Abschnitt~\ref{sec:results} fasst Benchmarks, Contests und Softwarequalit\"at
zusammen, Abschnitt~\ref{sec:conclusion} schlie\ss t mit einem Fazit.

\paragraph{Meilenstein\"ubersicht.}
\begin{center}
\begin{tabular}{clp{7.6cm}}
\toprule
\textbf{M} & \textbf{Thema} & \textbf{Kern} \\
\midrule
M1 & Dummy-KI      & Zuggenerator, Spielende-Erkennung, Random-Move \\
M2 & Basis-KI      & Alpha-Beta, iterative Tiefensuche, erste Bewertungsfunktion \\
M3 & Erweiterte KI & Transposition Table, PVS, Killer/History, Spielphasen \\
M4 & Optimierte KI & Contest-Anbindung, verbesserte Bewertung, MCTS (komplement\"ar) \\
\bottomrule
\end{tabular}
\end{center}

Entsprechend der Vorgabe wurde \textbf{kein Multithreading} eingesetzt; alle
Verfahren laufen single-threaded, um im Contest faire und reproduzierbare
Bedingungen zu gew\"ahrleisten.

% =====================================================================
\section{Das Spiel Tablut}\label{sec:tablut}
% =====================================================================
Tablut ist ein asymmetrisches Zwei-Parteien-Brettspiel auf einem
$9\times 9$-Feld.
Die beiden Parteien verfolgen unterschiedliche Ziele, was die KI-Entwicklung
besonders reizvoll macht.

\begin{itemize}[noitemsep]
  \item \textbf{Schwarz} (Angreifer) beginnt und stellt 16~Figuren.
        Ziel: den wei\ss en K\"onig fangen.
  \item \textbf{Wei\ss} (Verteidiger) stellt 8~Figuren plus den K\"onig.
        Ziel: den K\"onig in eines der vier Eckfelder fl\"uchten.
\end{itemize}

\paragraph{Zug- und Schlagregeln.}
\begin{itemize}[noitemsep]
  \item Alle Figuren ziehen wie T\"urme im Schach: orthogonal, beliebig weit,
        ohne \"uber Figuren zu springen.
  \item Die vier Eckfelder und das zentrale Thronfeld d\"urfen nur vom K\"onig
        betreten werden.
        Der leere Thron darf \"ubersprungen werden; hat der K\"onig den Thron
        verlassen, darf er nicht zur\"uckkehren.
  \item \textbf{Schlagen} (kustodial): Eine gegnerische Figur wird geschlagen,
        wenn sie orthogonal zwischen zwei feindlichen Figuren -- oder zwischen
        einer feindlichen Figur und einem Eckfeld bzw.\ dem (leeren) Thron --
        eingeklemmt wird.
  \item \textbf{K\"onig schlagen}: je nach Position sind 2 Angreifer (frei),
        3 Angreifer (neben dem Thron) oder 4 Angreifer (auf dem Thron) n\"otig.
  \item \textbf{Sieg}: Wei\ss\ gewinnt, wenn der K\"onig ein Eckfeld erreicht;
        Schwarz gewinnt, wenn der K\"onig geschlagen wird.
  \item \textbf{Remis}: Stellungswiederholung, kein legaler Zug oder eine
        festgelegte Zahl von Z\"ugen ohne Schlag.
\end{itemize}

Die asymmetrischen Ziele bedeuten, dass die Bewertungsfunktion beide Seiten
unterschiedlich behandeln muss (Abschnitt~\ref{sec:eval}); die Suche selbst
arbeitet daher nicht im Negamax-Schema, sondern behandelt Wei\ss\ als
Maximizer und Schwarz als Minimizer.

% =====================================================================
\section{Zustandsrepr\"asentation und Architektur}\label{sec:arch}
% =====================================================================
Das Brett wird durch \textbf{Bitboards} dargestellt (\code{src/board.py}):
Drei Ganzzahlen \code{attackers}, \code{defenders} und \code{king} kodieren je
eine Figurenart, wobei Bit~$r\cdot 9 + c$ das Feld $(r,c)$ belegt.
Zusammengeh\"orige abgeleitete Bitboards wie \code{occupied} sowie vorgerechnete
Masken (\code{CORNERS\_BB}, \code{EDGE\_BB}, \code{THRONE\_BIT}) erlauben schnelle
Mengenoperationen und Z\"ahlungen \"uber \code{bit\_count()} (Popcount auf
C-Ebene).
Erg\"anzend bietet \code{Board} feldweise Zugriffe (\code{get}, \code{set},
\code{is\_empty}) f\"ur Code, der einzelne Felder betrachtet.

\paragraph{Board als Hash-Schl\"ussel.}
\code{Board} definiert \code{\_\_hash\_\_} und \code{\_\_eq\_\_} auf Basis der Felder
\[
(\text{attackers},\ \text{defenders},\ \text{king},\ \text{side\_to\_move},\ \text{king\_left\_throne}).
\]
Diese beschreiben eine Stellung kollisionsfrei, sodass Stellungen ohne
zus\"atzliches Zobrist-Hashing als Dictionary-Schl\"ussel dienen k\"onnen -- das
wird in der Transposition Table ausgenutzt (Abschnitt~\ref{sec:tt}).

\paragraph{Module (Auswahl).}
\begin{center}
\begin{tabular}{ll}
\toprule
\textbf{Datei} & \textbf{Aufgabe} \\
\midrule
\code{board.py}                & Bitboard-Repr\"asentation, Hashing, Zugriffe \\
\code{moves.py}, \code{rules.py} & Zuggenerierung, Schlag- und Endregeln \\
\code{evaluation.py}           & Bewertungsfunktion, Spielphasen \\
\code{alpha\_beta.py}, \code{pvs.py} & Suchkerne \\
\code{transposition\_table.py} & Transposition Table \\
\code{search\_state.py}        & Killer Moves + History \\
\code{time\_manager.py}        & Zeitbudget je Spielphase \\
\code{iterative\_deepening.py} & Rahmen der iterativen Tiefensuche \\
\code{ai.py}                   & KI-Klassen (RandomAI \dots\ AdvancedAI) \\
\code{gameserver.py}, \code{server\_rules.py} & Contest-Anbindung (M4) \\
\bottomrule
\end{tabular}
\end{center}

\paragraph{Klassen- und Architektur\"uberblick.}
Abbildung~\ref{fig:classes} zeigt die zentralen Klassen und ihr Zusammenspiel.
Alle KI-Klassen teilen sich die Schnittstelle \code{choose\_move(board)}, sodass
sie im Turnier, im Selfplay und in Tests austauschbar sind.
Die produktive \code{AdvancedAI} orchestriert den Alpha-Beta-Stack (iterative
Tiefensuche, PVS, Transposition Table, Zugsortierung, Zeitmanagement); die
\code{MCTSAI} ist eine gleichwertige, aber unabh\"angige Alternative.
Der \code{TablutServerClient} verbindet eine KI mit dem TU-Gameserver und nutzt
dabei die server-treuen Regeln aus \code{server\_rules.py}.

\begin{figure}[htbp]
\centering
\begin{tikzpicture}[
    font=\footnotesize,
    box/.style={draw, rounded corners=2pt, align=left, inner sep=4pt,
                fill=gray!6, minimum height=7mm},
    core/.style={box, fill=codekw!12, draw=codekw!60},
    line/.style={-{Latex[length=2mm]}, gray!70},
    dep/.style={-{Latex[length=2mm]}, dashed, gray!60},
  ]
  % Interface + KIs
  \node[core] (adv) {\textbf{AdvancedAI}\\\code{choose\_move()}};
  \node[core, right=14mm of adv] (mcts) {\textbf{MCTSAI}\\\code{choose\_move()}};
  \node[box, above=7mm of $(adv)!0.5!(mcts)$] (iface)
        {\emph{gemeinsame Schnittstelle} \code{choose\_move(board)}};

  % Suchbausteine unter AdvancedAI
  \node[box, below=8mm of adv] (id) {\code{iterative\_deepening}};
  \node[box, below=6mm of id] (pvs) {\code{pvs} / \code{alpha\_beta}};
  \node[box, left=8mm of pvs] (tt) {\code{TranspositionTable}};
  \node[box, below=6mm of pvs] (ss) {\code{SearchState}\\\scriptsize Killer + History};
  \node[box, right=8mm of pvs] (tm) {\code{TimeManager}};

  % MCTS Bausteine
  \node[box, below=8mm of mcts] (node) {\code{MCTSNode}\\\scriptsize UCT + RAVE};

  % gemeinsame Basis
  \node[box, below=6mm of ss] (eval) {\code{evaluation}\\\scriptsize 7-Term-Heuristik, Spielphasen};
  \node[box, left=10mm of eval] (board) {\code{Board}\\\scriptsize Bitboards, Hash};
  \node[box, right=10mm of eval] (moves) {\code{moves} / \code{rules}};

  % Contest
  \node[box, right=14mm of mcts] (client) {\textbf{TablutServerClient}};
  \node[box, below=8mm of client] (srv) {\code{server\_rules}\\\scriptsize server-treue Schlaglogik};

  \draw[line] (iface) -- (adv);
  \draw[line] (iface) -- (mcts);
  \draw[line] (adv) -- (id);
  \draw[line] (id) -- (pvs);
  \draw[line] (pvs) -- (tt);
  \draw[line] (pvs) -- (ss);
  \draw[line] (pvs) -- (tm);
  \draw[line] (mcts) -- (node);
  \draw[dep] (pvs) -- (eval);
  \draw[dep] (node) -- (eval);
  \draw[dep] (eval) -- (board);
  \draw[dep] (eval) -- (moves);
  \draw[dep] (node) to[bend left=10] (moves);
  \draw[line] (client) -- (adv);
  \draw[dep] (client) -- (srv);
  \draw[dep] (srv) -- (board);
\end{tikzpicture}
\caption{Zentrale Klassen von DeepKnight. Durchgezogen: \glqq nutzt/steuert\grqq{};
gestrichelt: \glqq h\"angt ab von\grqq{}. Beide KI-Klassen implementieren dieselbe
\code{choose\_move}-Schnittstelle und teilen sich \code{Board}, Zuggenerierung und
Bewertung.}
\label{fig:classes}
\end{figure}

% =====================================================================
\section{Suchverfahren}\label{sec:search}
% =====================================================================
Die Suche entwickelte sich \"uber die Meilensteine von einfachem Minimax bis zur
kombinierten KI \code{AdvancedAI}, die alle Optimierungen vereint (Konfiguration
[7] in den Benchmarks, Abschnitt~\ref{sec:results}).
Als Rahmen dient die \emph{iterative Tiefensuche}.

% ---------------------------------------------------------------------
\subsection{Iterative Tiefensuche (Iterative Deepening)}
Statt mit fester Tiefe zu suchen, wird die Suche schrittweise vertieft
(Tiefe $1,2,3,\dots$).
Die Schleife l\"auft, solange \code{time\_manager.time\_left()} wahr ist; nach
jeder vollst\"andig abgeschlossenen Tiefe wird das Ergebnis \"ubernommen.
Vorteile: Es steht jederzeit ein g\"ultiger Zug bereit (gefahrloser Abbruch),
und die Eintr\"age aus flacheren Tiefen (Transposition Table, Killer/History)
verbessern die Zugsortierung der n\"achsttieferen Suche.
Das Verfahren ist damit ideal f\"ur harte Zeitlimits im Turnier.

% ---------------------------------------------------------------------
\subsection{Minimax und Alpha-Beta (M2)}
Die Basis-KI aus Meilenstein~2 nutzt Alpha-Beta-Pruning~\cite{russell2021}: Teilb\"aume, die das
Ergebnis nicht mehr beeinflussen k\"onnen, werden abgeschnitten.
Da Wei\ss\ als Maximizer und Schwarz als Minimizer behandelt wird (Score stets
aus Sicht von Wei\ss), gibt es getrennte Max-/Min-Zweige statt eines
Negamax-Kerns.
Bereits in M2 wurde eine Basis-Zugsortierung (Schlagz\"uge und K\"onigsz\"uge
Richtung Ecke zuerst) eingesetzt, da die G\"ute der Sortierung die Zahl der
Cutoffs bestimmt.

% ---------------------------------------------------------------------
\subsection{Transposition Table (M3, komplex)}\label{sec:tt}
W\"ahrend der Suche k\"onnen identische Stellungen \"uber verschiedene Zugfolgen
erreicht werden (\emph{Transpositionen})~\cite{cpw}.
Die Transposition Table (\code{src/transposition\_table.py}) speichert bereits
untersuchte Stellungen, um Mehrfachberechnung zu vermeiden.
Wie in Abschnitt~\ref{sec:arch} beschrieben, wird das \code{Board}-Objekt direkt
als Dictionary-Schl\"ussel verwendet -- kein Zobrist-Hashing, keine Kollisionen.

Jeder Eintrag speichert \code{depth}, \code{score}, \code{flag} und
\code{best\_move}.
Zur Alpha-Beta-Korrektheit werden drei Knotentypen unterschieden:

\begin{center}
\begin{tabular}{lll}
\toprule
\textbf{Typ} & \textbf{Wert} & \textbf{Bedeutung} \\
\midrule
\code{EXACT}       & 0 & exakter Minimax-Wert \\
\code{LOWER\_BOUND} & 1 & untere Schranke (Beta-Cutoff: Wert $\geq$ score) \\
\code{UPPER\_BOUND} & 2 & obere Schranke (kein Zug schlug $\alpha$: Wert $\leq$ score) \\
\bottomrule
\end{tabular}
\end{center}

Ein Eintrag wird nur genutzt, wenn seine Tiefe ausreicht
(\code{entry.depth $\geq$ depth}), und beim Speichern nur \"uberschrieben, wenn
die neue Suche mindestens so tief war.
Der gespeicherte \code{best\_move} wird beim n\"achsten Lookup als erster Zug
probiert und erzeugt so zus\"atzliche Cutoffs.

% ---------------------------------------------------------------------
\subsection{Principal Variation Search (M3, komplex)}
PVS (\code{src/pvs.py})~\cite{cpw} erweitert Alpha-Beta: Der erste Zug wird -- dank guter
Sortierung als wahrscheinlich bester Zug -- mit vollem Fenster $[\alpha,\beta]$
gesucht, alle weiteren zun\"achst mit einem \emph{Nullfenster}.
\"Uberschreitet ein Zug das enge Fenster nicht, ist keine weitere Suche n\"otig.
Nur bei einem \emph{Fail-High} wird mit vollem Fenster nachgesucht (Re-Search).
Da nicht im Negamax-Schema gearbeitet wird, existieren zwei Varianten:
\code{\_pvs\_max} f\"ur Wei\ss\ (Nullfenster
$[\alpha,\min(\beta,\alpha+\texttt{NULL\_WINDOW})]$) und \code{\_pvs\_min} f\"ur
Schwarz (Nullfenster $[\max(\alpha,\beta-\texttt{NULL\_WINDOW}),\beta]$), mit
\code{NULL\_WINDOW = 1.0}.

% ---------------------------------------------------------------------
\subsection{Erweiterte Zugsortierung: Killer Moves + History (M3)}
Zwei Heuristiken (\code{src/search\_state.py}) verbessern die Sortierung
weiter~\cite{cpw}.
\emph{Killer Moves}: Z\"uge, die auf einer Tiefe einen Beta-Cutoff ausgel\"ost
haben, sind oft auch in anderen Stellungen derselben Tiefe stark; pro Tiefe
werden zwei solche Z\"uge gespeichert.
\emph{History-Heuristik}: Eine globale Tabelle vergibt erfolgreichen
Cutoff-Z\"ugen einen Bonus von $\text{depth}^2$ und dient bei der Sortierung als
Tiebreaker.
Die resultierende Reihenfolge (\code{ordered\_for\_search}):
(1)~TT-\code{best\_move}, (2)~Killer Moves dieser Tiefe, (3)~restliche Z\"uge in
Basis-Sortierung, darin absteigend nach History-Score.
Die \"offentliche \code{order\_moves}-Schnittstelle blieb r\"uckw\"artskompatibel.

% ---------------------------------------------------------------------
\subsection{Spielphasen-Zeitmanagement (M3)}
Statt jedem Zug dasselbe Budget zuzuweisen, variiert es je nach Spielphase.
Die Phase bestimmt \code{game\_phase(board)} (\code{src/evaluation.py}) anhand
der Steinzahl, der K\"onigsposition (Thron/Rand) und der Distanz des K\"onigs zur
n\"achsten Ecke.
\code{src/time\_manager.py} skaliert das Basisbudget:

\begin{center}
\begin{tabular}{lc}
\toprule
\textbf{Phase} & \textbf{Zeitfaktor} \\
\midrule
Er\"offnung (\code{opening}) & 0.7 \\
Mittelspiel (\code{midgame}) & 1.3 \\
Endspiel (\code{endgame})    & 0.9 \\
\bottomrule
\end{tabular}
\end{center}

Das Mittelspiel ist strategisch am wichtigsten und erh\"alt mehr Zeit;
Er\"offnung und Endspiel sind oft eindeutiger.

% ---------------------------------------------------------------------
\subsection{Integration in \code{AdvancedAI}}
\code{AdvancedAI} (\code{src/ai.py}) vereint alle Techniken.
Pro Zug bestimmt \code{choose\_move(board)} die Spielphase und das angepasste
Zeitbudget, setzt Transposition Table und Suchzustand zur\"uck und ruft
\code{iterative\_deepening} mit \code{pvs\_with\_count} auf.
Transposition Table und \code{SearchState} bleiben \"uber alle Tiefen einer
Entscheidung erhalten; die Suche startet mit dem engen Fenster
$[-\texttt{WIN\_SCORE}, +\texttt{WIN\_SCORE}]$.

% =====================================================================
\section{Entwicklung der Bewertungsfunktion}\label{sec:eval}
% =====================================================================
Die Bewertungsfunktion \code{evaluate(board)} (\code{src/evaluation.py}) liefert
einen Wert aus Sicht von Wei\ss: positiv~=~gut f\"ur Wei\ss, negativ~=~gut f\"ur
Schwarz.

\paragraph{Anfangsversion (M1/M2).}
Zu Beginn wurde ausschlie\ss lich das Material bewertet:
$\text{Material} = \text{Verteidiger} - \text{Angreifer}$.

\paragraph{Positionelle Merkmale (M2$\rightarrow$M3).}
Nach und nach kamen positionelle Terme hinzu.

\paragraph{Verbesserungen im Meilenstein~4.}
F\"ur M4 wurde die schw\"achere Wei\ss-Verteidigung gezielt gest\"arkt: zwei neue
Terme, \code{corner\_approach\_score} (belohnt freie Schritte des K\"onigs in
Richtung einer Ecke -- anders als \code{escape\_lane\_score} muss der Weg nicht
vollst\"andig frei sein, es z\"ahlt der Fortschritt) und
\code{defender\_shield\_score} (belohnt Verteidiger direkt neben dem K\"onig als
Schutz-Eskorte).
Zus\"atzlich wurde \code{king\_distance\_score} auf die \emph{n\"achstgelegene}
Ecke umgestellt und die Gewichte neu justiert.
\code{ATTACKER\_VALUE} wurde nach Messung bei $1.0$ belassen, da die
Materialbewertung stabil war; die neuen Verteidiger-Terme verbessern Wei\ss,
ohne die Angriffsst\"arke von Schwarz zu besch\"adigen.
Die finale Funktion ist eine gewichtete Summe von sieben Teilbewertungen:

\begin{center}
\begin{tabular}{lll}
\toprule
\textbf{Teilbewertung} & \textbf{Beschreibung} & \textbf{Gewicht} \\
\midrule
\code{material\_score}        & Verteidiger $-$ Angreifer               & 1.0 \\
\code{king\_distance\_score}   & N\"ahe des K\"onigs zur n\"achsten Ecke & 3.0 \\
\code{king\_mobility\_score}   & Anzahl freier K\"onigsz\"uge            & 0.3 \\
\code{escape\_lane\_score}     & freie Fluchtwege zu den Ecken           & 4.5 \\
\code{king\_danger\_score}     & Angreifer in K\"onigsn\"ahe (negativ)   & 3.0 \\
\code{corner\_approach\_score} & Fortschritt Richtung Ecke \emph{(M4 neu)} & 1.8 \\
\code{defender\_shield\_score} & Verteidiger-Eskorte am K\"onig \emph{(M4 neu)} & 0.35 \\
\bottomrule
\end{tabular}
\end{center}

\paragraph{Contest-getriebene H\"artung (11.\,Juli).}
In den finalen Contests trat ein reproduzierbarer Schwachpunkt zutage: Als Wei\ss{}
lief der K\"onig deterministisch in Richtung einer Ecke, die bereits von Angreifern
bewacht war, und wurde dort geklemmt (K\"onig bei $(2,0)$ geschlagen -- in zwei
identischen Partien).
Zwei gezielte Erg\"anzungen beheben das:
\begin{itemize}
\item \textbf{Corner-Safety.} Die Anziehung jeder Ecke in \code{king\_distance\_score}
  und \code{corner\_approach\_score} wird mit der Zahl der Angreifer in Eck-N\"ahe
  (Manhattan-Distanz~$\le 4$) herunterskaliert. Eine bewachte Ecke zieht den K\"onig
  nicht mehr wie eine freie an.
\item \textbf{Klemm-Bedrohung.} \code{king\_danger\_score} bestraft nun zus\"atzlich
  (\code{KING\_THREAT\_WEIGHT}~$=8.0$) jede Stellung, in der ein Angreifer ein leeres
  K\"onig-Nachbarfeld in \emph{einem} Zug erreichen kann, w\"ahrend gegen\"uber bereits
  eine Klemm-Wand (Angreifer oder Eckfeld) steht -- also eine akute, sonst
  ``diagonal unsichtbare'' Schlagdrohung.
\end{itemize}
Nach dem Fix weicht die Suche der Verlustlinie fr\"uh aus (Selfplay: Remis statt
Wei\ss-Verlust in 17~Z\"ugen), und in den anschlie\ss enden Contests gewann die KI in
beiden Farbrollen (Abschnitt~\ref{sec:results}).

\paragraph{Terminale Stellungen und tiefenabh\"angiger Sieg.}
Terminale Stellungen erhalten den Betrag \code{WIN\_SCORE = 10000}.
Auf einen Hinweis aus der M2-R\"uckmeldung des Betreuers hin wurde der Sieg-Score
\emph{tiefenabh\"angig} gemacht, damit die KI den \emph{schnellsten} Sieg w\"ahlt
(\code{terminal\_score}):
\[
\text{Wei\ss-Sieg} = \texttt{WIN\_SCORE} - \text{plies\_from\_root},\qquad
\text{Schwarz-Sieg} = -\texttt{WIN\_SCORE} + \text{plies\_from\_root}.
\]

\paragraph{Performance-Refactoring.}
Die Effizienz der Bewertung wurde deutlich gesteigert: \code{king\_mobility\_score}
z\"ahlt K\"onigsz\"uge nun direkt auf den Bitboards statt \"uber
\code{generate\_legal\_moves}, und \code{king\_danger\_score} nutzt vorgerechnete
Nachbarschaftsmasken (\code{\_ADJACENT\_BB}, \code{\_NEAR\_BB}) mit
\code{bit\_count()}.
Gemessen auf der Startstellung: rund \textbf{215\,000 Bewertungen/s}
(vorher ca.\ 3\,000/s, also etwa $72\times$ schneller).

% =====================================================================
\section{Contest-Anbindung an den TU-Gameserver}\label{sec:contest}
% =====================================================================
F\"ur die Turnierteilnahme wurde im vierten Meilenstein ein eigener TCP-Client
f\"ur den offiziellen TU-Gameserver (gspy / Gameserver25) entwickelt
(\code{src/gameserver.py}).

\paragraph{Protokoll.}
Die Kommunikation ist zeilenbasiert \"uber TCP.
Der Client beherrscht Handshake und Login (\code{gspy}, \code{register},
\code{login}), das Erstellen bzw.\ Beitreten einer Lobby, die Konfigurationsphase
sowie die eigentliche Spielschleife.
Der Server \"ubertr\"agt nicht das komplette Brett, sondern nur die einzelnen
Z\"uge (Format \code{r,c,r,c}, 0-indiziert -- identisch zu unserem internen
Zug-Tupel).
Der Client pflegt daher ein eigenes \code{Board} und wendet jeden Zug -- eigenen
wie gegnerischen -- lokal an.

\paragraph{Server-treue Schlagregeln.}
Der Gameserver verwendet am Thron abweichende Schlagregeln.
Damit das lokal gepflegte Brett stets exakt dem Server-Brett entspricht (jede
Abweichung f\"uhrt zu einem illegalen Zug und damit zur Niederlage), wurde die
Server-Schlaglogik in \code{src/server\_rules.py} 1:1 nachgebildet
(\code{server\_captures}, \code{apply\_move\_server\_faithful}), inklusive der
Sonderf\"alle ``leerer Thron als Klemm-Wand'', ``Thron schl\"agt Schwarz immer''
und ``Eckfeld als Klemm-Wand''.
Diese Logik wird ausschlie\ss lich vom Client benutzt, nicht von der Suche.
\textbf{Verifikation:} \"uber $1\,000$ Zufallspartien
($152\,000$ Z\"uge) stimmten unser Brett und der echte Servercode zu
$100\,\%$ \"uberein -- \textbf{0 Desyncs, 0 illegale Z\"uge}.

\paragraph{Dynamisches, phasenabh\"angiges Zeitbudget.}
Der Client verwaltet das vom Server gemeldete Restzeit-Konto und bemisst das
Denkbudget pro Zug (\code{\_budget\_for\_move}) aus drei Faktoren: einem neutralen
Budget $\min(\text{base},\ \text{Restkonto}/20)$, das mit der Restzeit skaliert,
einem Spielphasen-Faktor (Er\"offnung $0{,}6$, Mittelspiel $1{,}5$, Endspiel $1{,}2$)
und harten Deckeln (nie unter $0{,}2$\,s, nie \"uber $1/6$ des Kontos bzw.\ das
Dreifache der Basis).
So denkt die KI im entscheidenden Mittelspiel am l\"angsten und verk\"urzt in der
Er\"offnung; da das Budget mit dem Restkonto schrumpft, verliert sie nie durch
Zeit\"uberschreitung.

\paragraph{Planbare Farbzuteilung.}
Der Gameserver ordnet die Farben \"uber den Lobby-Scheduler zu; der Default
\code{random} w\"urfelt sie pro Partie aus.
Damit die Seiten f\"ur die Auswertung planbar sind, setzt der Host-Client
\code{scheduler = round\_robin}: der Lobby-Ersteller wird Spieler~0 (Schwarz,
Angreifer), der Beitretende Spieler~1 (Wei\ss, Verteidiger).

\paragraph{Wiederholungs-Vermeidung.}
Da der Server kein automatisches Remis kennt, w\"urde stures Wiederholen einer
Stellung ins Leere laufen.
Der Client erkennt Stellungswiederholungen (\code{\_record\_position}) und weicht
nur dann auf einen alternativen Zug aus, wenn wir vorne liegen
(\code{\_repetition\_alternative}); andernfalls wird die Wiederholung
akzeptiert, statt einen erzwungenen Verlustzug zu spielen.

\paragraph{Live-Beobachtung.}
\code{watch\_board.py} zeigt die laufende Partie live an, inklusive
Gewonnen-/Verloren-/Unentschieden-Anzeige.

% =====================================================================
\section{Komplement\"are Technik: Monte-Carlo-Baumsuche}\label{sec:mcts}
% =====================================================================
Als komplement\"are Technik zum deterministischen Alpha-Beta-Stack wurde die
\textbf{Monte-Carlo-Baumsuche (MCTS)} mit dem UCT-Kriterium~\cite{russell2021}
umgesetzt (\"Uberblick in~\cite{browne2012}).
MCTS baut den Suchbaum stochastisch auf und ben\"otigt keine Trainingspipeline,
was gut zur symbolischen Ausrichtung des Projekts passt.
Jede Iteration durchl\"auft vier Phasen:

\begin{enumerate}[noitemsep]
  \item \textbf{Selection:} Vom Wurzelknoten wird nach der UCB1-Formel
        $\bar{x}_i + c\,\sqrt{\ln N / n_i}$ der vielversprechendste Ast
        abgestiegen.
  \item \textbf{Expansion:} Am Blatt wird ein neuer Kindknoten (Zug) angeh\"angt.
  \item \textbf{Simulation:} Von dort wird mit Zufallsz\"ugen bis zum Spielende
        ausgespielt (Rollout).
  \item \textbf{Backpropagation:} Das Ergebnis wird entlang des Pfades
        zur\"uckgez\"ahlt.
\end{enumerate}

Die Umsetzung nutzt die bestehende Zuggenerierung und Endregeln
(\code{apply\_move}, \code{is\_game\_over}) und ist -- wie der \"ubrige Code --
single-threaded.
Die MCTS-KI wird als \code{MCTSAI} mit derselben \code{choose\_move}-Schnittstelle
wie \code{AdvancedAI} bereitgestellt, sodass beide direkt gegeneinander antreten
k\"onnen.

\paragraph{Erweiterung: RAVE (Rapid Action Value Estimation)~\cite{browne2012}.}
Reines UCT hat unter knappem Zeitbudget ein Stichprobenproblem: Jeder Zug
sammelt nur \emph{eigene} Besuche, und bei einigen hundert Iterationen pro Zug
-- mehr sind in Python nicht drin -- bleibt die Statistik der meisten Z\"uge
praktisch leer.
\textbf{RAVE} begegnet dem mit der \emph{All-Moves-As-First}-Annahme (AMAF):
War ein Zug \emph{irgendwo sp\"ater} in einer Simulation gut, ist er
vermutlich auch \emph{jetzt} gut.
Jeder Knoten f\"uhrt daher neben den echten Besuchen eine zweite,
seitengetrennte AMAF-Statistik.
Der Rollout gibt seine gespielte Zugfolge zur\"uck; die Backpropagation
schreibt das Simulationsergebnis jedem Kindknoten gut, dessen Zug in dieser
Folge noch vorkommt -- auch wenn er nie ausgew\"ahlt wurde.
So entstehen pro Iteration statt eines einzigen Datenpunkts Dutzende.

Beide Sch\"atzer werden in der Selektion gemischt.
Der RAVE-Wert ist breit, aber ungenau; der UCT-Wert ist exakt, aber langsam.
Das Gewicht
\[
  \beta = \sqrt{\frac{k}{3n + k}}, \qquad k = 300,
\]
sorgt f\"ur den \"Ubergang: Bei $n = 0$ echten Besuchen dominiert RAVE, mit
wachsendem $n$ f\"allt $\beta$ gegen null und die exakte UCT-Statistik
\"ubernimmt.
\"Uber den Schalter \code{use\_rave} l\"asst sich die Erweiterung f\"ur
A/B-Messungen abschalten.

\paragraph{Spielst\"arke-Messung (Selfplay).}
Ein leichtgewichtiger, deterministischer Selfplay-Helfer
(\code{python -m src.m4\_selfplay}) l\"asst \code{MCTSAI} und \code{AdvancedAI}
in beiden Farbrollen gegeneinander antreten
(Seed~7, $0{,}2$\,s pro Zug, \code{AdvancedAI} bis Tiefe~2, MCTS mit
$120$~Iterationen).
Ergebnis: In beiden Partien gewann die \emph{schwarze} (Angreifer-)Seite --
\code{AdvancedAI} als Schwarz in 55~Z\"ugen, \code{MCTSAI} als Schwarz in
37~Z\"ugen, also \textbf{1:1}.
Das best\"atigt einerseits den strukturellen Vorteil der Angreifer im Tablut und
zeigt andererseits, dass MCTS unter diesem knappen Budget durchaus
konkurrenzf\"ahig ist.
Der ausgereifte Alpha-Beta-Stack bleibt jedoch die contest-erprobte Haupt-KI;
MCTS dient als komplement\"are Technik und Vergleichsbasis.

\paragraph{Wirkung von RAVE -- ein negativer Befund.}
Um den Nutzen von RAVE zu belegen statt ihn nur zu behaupten, traten beide
Varianten bei identischem Iterationsbudget (200~Iterationen pro Zug) direkt
gegeneinander an, mit wechselnden Farbrollen.
Das Ergebnis liegt mit 7:5 aus Sicht der RAVE-Variante im Rauschen: Bei dieser
Stichprobengr\"o\ss e w\"are ein solcher Ausgang auch zwischen zwei identischen
Gegnern unauff\"allig.
Eine zweite Messung \"uber 30~zuf\"allig angespielte Stellungen mit 150 bzw.\
400~Iterationen zeigte ebenfalls keinen belastbaren Unterschied.
Wir k\"onnen f\"ur RAVE in unserem Aufbau also \emph{keine} Verbesserung der
Spielst\"arke nachweisen.

Daf\"ur gibt es zwei plausible Gr\"unde.
Erstens ist unser Iterationsbudget zu klein: RAVE stammt aus Go-Programmen mit
vielen tausend Playouts pro Zug, w\"ahrend unsere Python-Implementierung im
Contest-Zeitrahmen nur einige hundert schafft -- der Baum bleibt so flach, dass
sich der Informationsvorsprung der AMAF-Statistik kaum auszahlen kann.
Zweitens verletzt Tablut die AMAF-Grundannahme st\"arker als Go: Ob ein Zug gut
ist, h\"angt hier massiv von der K\"onigsstellung ab, sodass derselbe Zug je nach
Zeitpunkt hervorragend oder verlustbringend sein kann.
Wir dokumentieren dies bewusst als Ergebnis: Eine gemessene Nullaussage ist
aussagekr\"aftiger als eine unbelegte Verbesserung, und die Implementierung
bleibt \"uber \code{use\_rave} f\"ur sp\"atere Messungen mit gr\"o\ss erem
Budget erhalten.

% =====================================================================
\section{Benchmarks, Contests und Softwarequalit\"at}\label{sec:results}
% =====================================================================

\subsection{Suchbenchmarks}
Die Benchmarks wurden mit \code{python -m src.m3\_experiments} bei $1.0$\,s
Zeitlimit pro Konfiguration durchgef\"uhrt (Apple~M1~Pro, macOS, Python~3.13.2).
Verglichen werden sieben Konfigurationen von reinem Alpha-Beta bis zu
\code{AdvancedAI}.
Alle liefern denselben besten Zug und Score -- die Optimierungen \"andern nur die
Effizienz, nicht das Ergebnis.
Tabelle~\ref{tab:nodes} zeigt f\"ur drei repr\"asentative Stellungen die erreichte
Tiefe~(T) und die untersuchten Knoten~(K).

\begin{center}
\captionof{table}{Erreichte Tiefe (T) und untersuchte Knoten (K) bei 1\,s}\label{tab:nodes}
\small
\setlength{\tabcolsep}{4.5pt}
\begin{tabular}{>{\raggedright}p{5.3cm}cccccc}
\toprule
 & \multicolumn{2}{c}{\textbf{Start}} & \multicolumn{2}{c}{\textbf{Wei\ss\ in 2}} & \multicolumn{2}{c}{\textbf{Endspiel}} \\
\textbf{Konfiguration} & T & K & T & K & T & K \\
\midrule
{[1]} Alpha-Beta (M2-Baseline)     & 3 & 5781 & 5 & 3710 & 5 & 2836 \\
{[2]} AB + TT                      & 3 & 8642 & 6 & 2360 & 6 & 2463 \\
{[3]} AB + Killer + History       & 3 & 7093 & 5 & 3383 & 7 & 2943 \\
{[4]} PVS                         & 3 & 5857 & 5 & 3827 & 5 & 2836 \\
{[5]} PVS + TT                    & 3 & 6369 & 6 & 2353 & 6 & 2412 \\
{[6]} PVS + TT + Killer + History & 3 & 7498 & 6 & \textbf{2005} & 6 & \textbf{2073} \\
{[7]} AdvancedAI                   & 3 & 7498 & \textbf{6} & \textbf{2005} & \textbf{6} & \textbf{2073} \\
\bottomrule
\end{tabular}
\end{center}

Der Gewinn zeigt sich zweifach: gr\"o\ss ere erreichte Tiefe im gleichen Budget
und weniger Knoten bei gleicher Tiefe.
In ``Wei\ss\ gewinnt in 2'' und im Endspiel erreicht \code{AdvancedAI} bei $1$\,s
Tiefe~6, w\"ahrend reines Alpha-Beta nur Tiefe~5 schafft.
Die Transposition Table wirkt am st\"arksten in transpositionsreichen
Endspielstellungen, die erweiterte Zugsortierung vor allem im Mittelspiel.

\subsection{Contests}
Die Spielst\"arke wurde in echten Contests gegen eine andere Kursgruppe gemessen;
die Verbindung lief \"uber einen \code{bore}- bzw.\ \code{ngrok}-TCP-Tunnel, mit
$5$\,s Basis-Bedenkzeit pro Zug und $120$\,s Zeitkonto pro Spieler.
Aus fr\"uheren Contests wurden zwei Ursachen behoben: ein zu knapp gesetztes
\code{time\_account} (Gegner-Flag/-Crash) und Endlos-Wiederholungen
(Abschnitt~\ref{sec:contest}).
Die finalen Partien am 11.\,Juli deckten zudem die K\"onig-Falle auf
(Abschnitt~\ref{sec:eval}): Als Wei\ss{} verlor die KI zun\"achst, weil der K\"onig in
eine bewachte Ecke lief.
Nach der Corner-Safety-/Klemm-Bedrohungs-H\"artung gewann sie die Folgepartien in
\emph{beiden} Farbrollen.

\begin{center}
\begin{minipage}{\linewidth}
\centering
\captionof{table}{Finale Contests gegen eine andere Kursgruppe (11.\,Juli 2026)}
\begin{tabular}{ll>{\raggedright\arraybackslash}p{5.6cm}}
\toprule
\textbf{Partie} & \textbf{Farbe} & \textbf{Ergebnis} \\
\midrule
vor Eval-H\"artung  & Wei\ss   & Niederlage -- K\"onig bei $(2,0)$ geklemmt ($17$~Halbz\"uge) \\
nach Eval-H\"artung & Wei\ss   & Sieg -- K\"onig-Flucht in die Ecke ($18$--$22$~Halbz\"uge) \\
nach Eval-H\"artung & Schwarz  & Sieg -- K\"onig gefangen ($41$~Halbz\"uge) \\
\bottomrule
\end{tabular}
\end{minipage}
\end{center}

\subsection{Softwarequalit\"at und Tests}
Die Korrektheit ist durch eine umfangreiche Testsuite abgesichert:
\textbf{138 Tests}, alle gr\"un in der finalen Projektfassung.
Getestet werden u.\,a.\ PVS, Transposition Table, Zugsortierung/Suchzustand,
Zeitmanagement, die tiefenabh\"angigen Sieg-Scores, die kombinierte KI sowie die
MCTS-Suche (\code{test\_mcts.py}, \code{test\_m4\_selfplay.py}).
Insbesondere wird gepr\"uft, dass PVS bei gleicher Tiefe dasselbe Ergebnis
liefert wie Alpha-Beta.
Der gesamte Code liegt versioniert im Repository
\href{https://github.com/BatuCani/deepknight}{github.com/BatuCani/deepknight}.

% =====================================================================
\section{Entwicklung und Erkenntnisse}\label{sec:process}
% =====================================================================
Die Funktionalit\"at wurde inkrementell aufgebaut: Grundger\"ust,
Basissuche, erweiterte Suchverfahren sowie Optimierung und Serveranbindung.

\paragraph{Werkzeuge.}
Der Quellcode wird mit \textbf{Git} versioniert.
Als Sprache dient \textbf{Python~3} ausschlie\ss lich mit der Standardbibliothek;
einzige externe Abh\"angigkeit ist \textbf{pytest} f\"ur die Tests.
Der Bericht ist in \LaTeX{} gesetzt.

\paragraph{Qualit\"atssicherung.}
Zentrale Ma\ss nahme ist eine wachsende \textbf{Testsuite},
die Regressionen in Zuggenerierung, Regeln und Suchverfahren sichtbar macht.
Kritische Invarianten sind explizit abgesichert -- etwa dass PVS bei gleicher
Tiefe dasselbe Ergebnis liefert wie reines Alpha-Beta, oder dass die
server-treue Schlaglogik \"uber $152\,000$ Z\"uge nicht vom echten Servercode
abweicht (0~Desyncs).
So wurden Korrektheit und Turnierfestigkeit messbar statt nur behauptet.

\paragraph{Reflexion.}
R\"uckblickend war die gr\"o\ss te Lehre, dass reale Contests Schwachstellen
aufdecken, die weder Tests noch Selfplay zeigten: Erst eine reproduzierbare
Niederlage als Wei\ss{} f\"uhrte zur \glqq K\"onig-Falle\grqq{}-H\"artung
(Abschnitt~\ref{sec:eval}).
Ebenso zahlte sich der fr\"uhe Umstieg auf Bitboards aus -- der Geschwindigkeits\-gewinn
(rund $72\times$) erm\"oglichte tiefere Suchen im selben Zeitbudget.
Nicht jede Erweiterung brachte messbaren Gewinn: F\"ur RAVE lie\ss{} sich in
unserem Aufbau keine St\"arkeverbesserung nachweisen -- ein bewusst dokumentierter
negativer Befund (Abschnitt~\ref{sec:mcts}).

% =====================================================================
\section{Fazit und Ausblick}\label{sec:conclusion}
% =====================================================================
\"Uber vier Meilensteine ist mit \code{DeepKnight} ein vollst\"andiger,
turnierf\"ahiger Alpha-Beta-Stack entstanden: Bitboard-Repr\"asentation,
Transposition Table, Principal Variation Search, erweiterte Zugsortierung und
Spielphasen-Zeitmanagement, kombiniert mit einer positionell reichen
Bewertungsfunktion.
Im vierten Meilenstein kamen die Contest-Anbindung an den TU-Gameserver mit
server-treuer Schlaglogik (verifiziert \"uber $152\,000$ Z\"uge), eine gest\"arkte
Wei\ss-Verteidigung sowie die komplement\"are Monte-Carlo-Baumsuche hinzu.
Die finale Projektfassung ist durch 138~Tests abgesichert.

\paragraph{Ausblick.}
Weiteres Potenzial liegt in der Feinjustierung der Bewertungsgewichte (etwa per
evolution\"arem Lernen), einer Ruhesuche zur Vermeidung des Horizont-Effekts
sowie einer engeren Verzahnung von MCTS und Alpha-Beta (z.\,B.\ MCTS mit
Alpha-Beta-Rollouts statt reiner Zufallsz\"uge).

% =====================================================================
\begin{thebibliography}{9}
% =====================================================================
\bibitem{russell2021}
S.~Russell und P.~Norvig.
\emph{Artificial Intelligence: A Modern Approach.}
4.~Auflage, Pearson, 2021.
(Grundlagen der Spielbaumsuche: Minimax, Alpha-Beta, Monte-Carlo-Baumsuche/UCT.)

\bibitem{cpw}
\emph{Chessprogramming Wiki.}
\url{https://www.chessprogramming.org}, abgerufen im Juli 2026.
(Referenz f\"ur die praktischen Such-Erweiterungen: Transposition Table,
Principal Variation Search, Killer Moves, History-Heuristik.)

\bibitem{browne2012}
C.~Browne u.\,a.
\emph{A Survey of Monte Carlo Tree Search Methods.}
IEEE Transactions on Computational Intelligence and AI in Games,
4(1):1--43, 2012.
(\"Ubersicht zu MCTS-Varianten einschlie\ss lich AMAF/RAVE.)
\end{thebibliography}

\end{document}
